%%%%%%%%%%%%%%%%%%%%%%%%%%%%%%%%%%%%%%%%%%%%%%%%%%%%%%%%%%%%%%%%
%%                   Semester project 1	  	              %%
%%              Smart home energy monitoring   		      %%
%%                       Andrew Watson                        %%
%%                           MT-MA1                           %%
%%%%%%%%%%%%%%%%%%%%%%%%%%%%%%%%%%%%%%%%%%%%%%%%%%%%%%%%%%%%%%%%

%\documentclass[a4paper,draft]{article}
\documentclass[a4paper]{article}
% NB : enlever ``draft'' pour avoir les images!

\usepackage[english]{babel}
\usepackage[utf8]{inputenc}
\usepackage[T1]{fontenc} % for hyphenation of accentuated words
\usepackage{ae,aecompl}  % alternative to the fugly T1 fonts
\usepackage[pdftex]{graphicx} % modern graphics package
\usepackage{multirow}   % multiple rows in tables

%\usepackage{microtype}
%\usepackage{hyph-utf8}

\usepackage{wasysym}    % extended symbols
\usepackage{includes/mystyle}    % added styling
\usepackage{marvosym}   % euro symbol, among other things.
\usepackage[eurosym]{eurofont}

\usepackage[usenames,dvipsnames]{xcolor}  % if using PDFLaTeX
%\usepackage[usenames,dvips]{color}      % for regular LaTeX 

\usepackage{fancyhdr}   % fancy headers
\pagestyle{fancy}

%\usepackage{textcomp}  % some additionnal symbols
\usepackage{units}      % pretty units : \unit[val]{dim}
\usepackage{url}        % pretty urls
\usepackage{listings}   % pretty code listings \begin{lstlisting}
\usepackage{tikz} % pic/graph package. used for rounded corners :)

\graphicspath{{.}{./images/}}
\usepackage{subfigure} % multiple figures

%\addtocounter{secnumdepth}{1}  % number paragraphs

\usepackage{lscape}     % use landscape mode with \begin{landscape}
\usepackage{wrapfig}    % invoke Satan and wrap figures in text

\usepackage{hyperref}

\hypersetup{
    bookmarks=true,         % show bookmarks bar?
    unicode=true,          % non-Latin characters in Acrobat’s bookmarks
    pdftoolbar=true,        % show Acrobat’s toolbar?
    pdfmenubar=true,        % show Acrobat’s menu?
    pdffitwindow=true,     % window fit to page when opened
    pdfstartview={FitV},    % fits the width of the page to the window
    pdftitle={1st Semester Thesis},    % title
    pdfauthor={Andrew Watson},     % author
    pdfsubject={Wireless smart home network},   % subject of the document
    pdfcreator={Andrew Watson},   % creator of the document
    pdfproducer={EPFL}, % producer of the document
    pdfkeywords={AVR, embedded, ZigBee, Powerline, ELab}, % list of keywords
    pdfnewwindow=true,      % links in new window
    colorlinks=true,       % false: boxed links; true: colored links
    linkcolor=MidnightBlue,          % color of internal links
    citecolor=RoyalPurple,        % color of links to bibliography
    filecolor=magenta,      % color of file links
    urlcolor=DarkOrchid           % color of external links
}

\begin{document}
\begin{titlepage}
\nocite{*}      % to make sure bibliography appears in the correct order
  \begin{center}
     
     
    % Upper part of the page
    %\includegraphics[width=4cm]{logo_epfl}\\[0.5cm]
     
    \vfill 
    \textsc{\LARGE Ecole Polytechnique Fédérale de Lausanne }\\[1.0cm]

    { \huge \bfseries Smart home energy monitoring}\\[0.4cm]
    \includegraphics[width=12cm]{images/logo_epfl}\\[0.5cm]

    \vfill

     
    % Bottom of the page
    Andrew \textsc{Watson}\\[0.5cm] 
    Version 0.1.1\\[0.5cm]
    {\large \today} 
     
  \end{center}

\end{titlepage}
%\maketitle

\newpage{}

%\fancyhead{}
\fancyfoot{}
\lhead{}
\cfoot{\thepage}        % numéro de page..
\lfoot{Semester Project}
%\rfoot{\today}
\rfoot{\today} %% TODO : fix the date

%\begin{abstract}
%\end{abstract}

% CUSTOM COMMANDS FOR THIS REPORT
\newcommand{\expr}[1]{\og \emph{{#1}} \fg} % \expr{word} => « Word »
\newcommand{\important}[1]{\textbf{#1}}
\newcommand{\code}[1]{\texttt{#1}}

%\renewcommand\contentsname{Plan}  % Rename ``Table des Matières''
\tableofcontents{}

\newpage

%%%% BEGIN : LSTLISTINGS CONFIG %%%%%%
%%%% TODO : MOVE TO SEPARATE FILE ONCE FINISHED %%%%%
%% see http://www.jorgemarsal.com/blog/2009/06/08/source-code-snippets-in-latex/
\lstset{language=C}
%\definecolor{lightgrey}{RGB}{200,200,200}
\definecolor{grey92}{gray}{0.92}
\definecolor{grey75}{gray}{0.75}
\definecolor{grey45}{gray}{0.45}
\lstdefinestyle{console}
{
  numbers=none,
  basicstyle=\bf\ttfamily,
  backgroundcolor=\color{grey92},
  frame=lrtb,
  framerule=0.5pt,
  linewidth=\textwidth,
}
\lstdefinestyle{avr-c}
{
  style=console
}

\lstset{
  style=console
}

%%%%%%% END : LSTLISTINGS CONFIG %%%%%%%%


\section*{Introduction}
\addcontentsline{toc}{section}{Introduction}
\markboth{Introduction}{\MakeUppercase{Introduction}}
This is the introduction for my semester project

\section{Some examples of code}
\begin{lstlisting}[style=avr-c]
  int main(void)
  {
    code_example();
  }
\end{lstlisting}

There's also a way of including \code{code} inline.

\end{document}
